% Generated draft: requires mathematical and source review.
\documentclass[11pt]{article}
\usepackage[margin=25mm]{geometry}
\usepackage{fontspec}
\setmainfont{Latin Modern Roman}
\usepackage{amsmath,amssymb,mathtools,unicode-math}
\setmathfont{Latin Modern Math}
\usepackage{enumitem,needspace,xurl,hyperref}
\setlength{\parindent}{0pt}
\setlength{\parskip}{5pt}
\setlength{\emergencystretch}{3em}
\newcommand{\N}{\mathbb{N}}
\newcommand{\Z}{\mathbb{Z}}
\newcommand{\Q}{\mathbb{Q}}
\newcommand{\R}{\mathbb{R}}
\newcommand{\C}{\mathbb{C}}
\newcommand{\F}{\mathbb{F}}
\newcommand{\A}{\mathbb{A}}
\newcommand{\PP}{\mathbb{P}}
\newcommand{\E}{\mathbb{E}}
\newcommand{\eps}{\varepsilon}
\newcommand{\dd}{\,\mathrm{d}}
\begin{document}
\section*{UnsolvedMath ID: 30005090}
\textbf{Source problem code:} catalogue entry
\subsection{Definitions and mathematical statement}
A regular tetrahedron in \(\R^3\) is the convex hull of four affinely independent points whose six pairwise distances are equal to some positive number. Its vertices are those four points, and \(\operatorname{int}(T)\) denotes its interior in \(\R^3\).

For a positive integer \(n\), call a family \(T_1,\ldots,T_n\) admissible if each \(T_i\) is a regular tetrahedron, there exists \(p\in\R^3\) that is a vertex of every \(T_i\), and
\[
\operatorname{int}(T_i)\cap\operatorname{int}(T_j)=\varnothing
\qquad(1\leq i<j\leq n).
\]
Thus boundary contacts, including contacts beyond the common vertex, are permitted. The supplied statement does not explicitly require different tetrahedra to have equal edge lengths; this formulation imposes no such additional restriction.

Determine the maximum
\[
\tau_3=\max\bigl\{n\in\Z_{\geq1}:\text{there exists an admissible family of }n\text{ tetrahedra}\bigr\}.
\]
The target is the exact value of \(\tau_3\), with a construction attaining it and an upper bound excluding larger admissible families.

The supplied dated curation states that an icosahedral construction gives \(20\) tetrahedra and that the report improves the elementary upper bound \(22\), while neither the exact value nor the optimality of \(20\) is proved there. The improved numerical upper bound is not supplied. These are source-reported qualifications, not independently verified conclusions of this draft.
\subsection{Short English statement}
Determine the largest number \(\tau_3\) of regular tetrahedra in three-dimensional Euclidean space that can have one common vertex while their interiors remain pairwise disjoint. Boundary contacts are allowed. The supplied record labels the problem open; its dated assessment reports a construction with \(20\) tetrahedra without asserting that this number is optimal.
\subsection{Sources}
\begin{itemize}
\item Supplied local UnsolvedMath registry record, permanent ID \(30005090\), title ``Kissing Number of Regular Tetrahedra,'' in the collection ``Oberwolfach Reports Open Problems.'' Exact source problem code: \texttt{catalogue entry}.
\item Supplied report citation: \emph{Conic Linear Optimization for Computer-Assisted Proofs} (2023), Oberwolfach Reports, pp.~1039--1089, DOI \texttt{10.4171/owr/2022/20}. The original statement identifies p.~1079 and the contribution ``The kissing number problem for regular tetrahedra,'' by Andreas Spomer, joint work with Frank Vallentin, Fernando M\'{a}rio de Oliveira Filho, Stefan Krupp, and Fabr\'{i}cio C. Machado. The registry separately gives the proposed year as 2022. Supplied literature-source label: ``Spomer et al., regular-tetrahedron kissing number.'' \url{https://doi.org/10.4171/owr/2022/20}
\item Supplied literature source labeled ``Oberwolfach primary abstract'': \url{https://oa.tib.eu/renate/bitstreams/c99d51a5-c8ce-41bf-afb1-b59e6547c379/download}
\item Literature assessment checked at \texttt{2026-08-22}: supplied status ``open''; supplied assessment reports the icosahedral construction with \(20\) tetrahedra, an improvement to the elementary upper bound \(22\), and no proof of the exact value or optimality of \(20\). This is supplied dated curation derived from the Oberwolfach report collection, not independently verified literature assessment. The record requests independent verification of time-sensitive status and cited claims.
\item The registry labels the statement \texttt{corrected\_verified} and reports that it was checked against the official Oberwolfach report for preservation of hypotheses, notation, and target. This verification claim is supplied provenance; this draft does not claim to have visited or checked either external source.
\end{itemize}
\end{document}
